% ==========================================
% FRONT MATTER - PRELIMINARY PAGES
% ==========================================

% --- 1. TITLE PAGE ---
\begin{titlepage}
\begin{center}
    \setstretch{1.2}
    {\large \textbf{VISVESVARAYA TECHNOLOGICAL UNIVERSITY}} \\
    {\large \textbf{BELAGAVI - 590018, KARNATAKA}} \\
    \vspace{1.5cm}
    \begin{figure}[h!]
        \centering
        % \includegraphics[width=3cm]{vtu_logo.png} % Placeholder for VTU Logo
        \vspace{2cm} % Space instead of logo for compilation safety
    \end{figure}
    \vspace{0.5cm}
    {\Large \textbf{FEATURE-ENHANCED HANDWRITTEN TEXT RECOGNITION USING CNN-BiLSTM WITH ATTENTION AND NLP POST-PROCESSING}} \\
    \vspace{1.5cm}
    \textit{A thesis report submitted in partial fulfillment of the requirements for the award of degree of} \\
    \vspace{0.5cm}
    {\large \textbf{MASTER OF COMPUTER APPLICATIONS}} \\
    \vspace{1.0cm}
    \textit{Submitted by} \\
    \vspace{0.3cm}
    {\large \textbf{Nagaraj Aithal}} \\
    {\large \textbf{USN: 4MT24MCA000}} \\
    \vspace{1.0cm}
    \textit{Under the guidance of} \\
    \vspace{0.3cm}
    {\large \textbf{Dr. Madhwaraj K G}} \\
    \textit{Guide, Department of MCA, MITE} \\
    \vspace{1.5cm}
    % \includegraphics[width=2.5cm]{mite_logo.png} % Placeholder for MITE Logo
    \vspace{1.0cm}
    \textbf{DEPARTMENT OF MASTER OF COMPUTER APPLICATIONS} \\
    \textbf{MANGALORE INSTITUTE OF TECHNOLOGY \& ENGINEERING} \\
    \setstretch{1.0}
    {\small (An Autonomous Institution, Affiliated to VTU Belagavi, Approved by AICTE, Accredited by NAAC)} \\
    {\small Mijar, Moodabidri, Karnataka 574225} \\
    \vspace{0.8cm}
    \textbf{JULY 2026 -- OCTOBER 2026}
\end{center}
\end{titlepage}

\newpage

% --- 2. DEPARTMENT CERTIFICATE ---
\begin{center}
    \setstretch{1.3}
    \textbf{MANGALORE INSTITUTE OF TECHNOLOGY \& ENGINEERING} \\
    {\small (An Autonomous Institution, Affiliated to VTU Belagavi, Approved by AICTE, Accredited by NAAC)} \\
    \textbf{Mijar, Moodabidri, Karnataka 574225} \\
    \vspace{0.5cm}
    \textbf{DEPARTMENT OF MASTER OF COMPUTER APPLICATIONS} \\
    \vspace{1.5cm}
    \Large{\textbf{CERTIFICATE}} \\
    \vspace{1.0cm}
\end{center}
\setstretch{1.5}
This is to certify that the Major Project Phase - 2 (24MCSE621) work entitled \textbf{"Feature-Enhanced Handwritten Text Recognition Using CNN-BiLSTM with Attention and NLP Post-Processing"} is a bonafide work carried out by \textbf{Nagaraj Aithal} (USN: \textbf{4MT24MCA000}) during the academic year 2025-2026 in partial fulfillment of the requirements for the award of the degree of \textbf{MASTER OF COMPUTER APPLICATIONS} under Visvesvaraya Technological University, Belagavi, Karnataka. The thesis has been approved as it satisfies the academic requirements in respect of project work prescribed for the said degree.
\vspace{3cm}
\begin{table}[h!]
    \centering
    \setstretch{1.0}
    \begin{tabular}{ccc}
        \textbf{Dr. Madhwaraj K G} & \textbf{Dr. Madhwaraj K G} & \textbf{Dr. Prashanth C M} \\
        Guide / Supervisor & Head of the Department & Principal \\
        Department of MCA & Department of MCA & MITE \\
    \end{tabular}
\end{table}

\newpage

% --- 3. GUIDE CERTIFICATE ---
\begin{center}
    \setstretch{1.3}
    \Large{\textbf{CERTIFICATE}} \\
    \vspace{1.5cm}
\end{center}
\setstretch{1.5}
This is to certify that the work contained in the thesis entitled \textbf{"Feature-Enhanced Handwritten Text Recognition Using CNN-BiLSTM with Attention and NLP Post-Processing"}, submitted by \textbf{Nagaraj Aithal} (USN: \textbf{4MT24MCA000}) for the award of the degree of \textbf{Master of Computer Applications} to Visvesvaraya Technological University, Belagavi, is a record of bonafide research work carried out by him/her under my direct supervision and guidance.

I consider the thesis has reached the standards and fulfills the requirements of the rules and regulations related to the MCA degree. The contents embodied in the thesis have not been submitted for the award of any other degree in this College or any other College/University.

\vspace{4cm}
\begin{flushleft}
    Date: \\
    Place: Moodabidri
\end{flushleft}
\begin{flushright}
    \vspace{-1.5cm}
    \textbf{Signature of Supervisor} \\
    \textbf{Dr. Madhwaraj K G} \\
    Guide, Department of MCA \\
\end{flushright}

\newpage

% --- 4. DECLARATION ---
\begin{center}
    \setstretch{1.3}
    \Large{\textbf{DECLARATION}} \\
    \vspace{1.5cm}
\end{center}
\setstretch{1.5}
I certify that:
\begin{enumerate}
    \item[a.] The work contained in the thesis is original and has been done by myself under the supervision of my supervisor.
    \item[b.] The work has not been submitted to any other Institute for any degree or diploma.
    \item[c.] I have conformed to the norms and guidelines given in the Ethical Code of Conduct of the Institute.
    \item[d.] Whenever I have used materials (data, theoretical analysis, and text) from other sources, I have given due credit to them by citing them in the text of the thesis and giving their details in the references.
    \item[e.] Whenever I have quoted written materials from other sources, due credit has been given to the sources by citing them.
    \item[f.] From the plagiarism test, it is found that the similarity index of the whole thesis is within 8\% and that of any single paper is less than 10\%, in accordance with the university guidelines.
\end{enumerate}

\vspace{3cm}
\begin{flushleft}
    Date: \\
    Place: Moodabidri
\end{flushleft}
\begin{flushright}
    \vspace{-1.5cm}
    \textbf{Nagaraj Aithal} \\
    \textbf{USN: 4MT24MCA000} \\
\end{flushright}

\newpage

% --- 5. ACKNOWLEDGEMENTS ---
\begin{center}
    \Large{\textbf{ACKNOWLEDGEMENTS}} \\
    \vspace{1.0cm}
\end{center}
\setstretch{1.5}
I would like to express my sincere gratitude to everyone who guided and supported me throughout this thesis work.

My deepest thanks go to my guide, \textbf{Dr. Madhwaraj K G}, Department of Master of Computer Applications, MITE, for his constant encouragement, technical insights, and constructive feedback at every stage of this project. The clarity of direction provided by my guide made what seemed like a complex undertaking feel structured and achievable.

I am extremely grateful to \textbf{Dr. Madhwaraj K G}, Head of the Department, MCA, for providing a highly conducive research environment and for making departmental resources readily accessible. His leadership and support have been invaluable to the entire MCA program.

I extend my sincere thanks to \textbf{Dr. Prashanth C M}, Principal, Mangalore Institute of Technology \& Engineering, for his leadership and for fostering a culture of research and academic excellence within the institution.

The faculty members of the Department of Master of Computer Applications have been a consistent source of knowledge and motivation. I am thankful to each of them for the academic foundation they helped build over the course of this program.

I am also grateful to the technical staff and library team at MITE for their timely support in providing research materials and computing facilities.

Finally, I owe a great deal to my family and friends, whose patience, encouragement, and belief in me made this journey possible. This work is as much theirs as it is mine.

\vspace{2cm}
\begin{flushright}
    \textbf{Nagaraj Aithal} \\
    \textbf{USN: 4MT24MCA000}
\end{flushright}

\newpage

% --- 6. ABSTRACT ---
\begin{center}
    \Large{\textbf{ABSTRACT}} \\
    \vspace{1.0cm}
\end{center}
\setstretch{1.5}
Handwritten Text Recognition (HTR) is a technically demanding problem in document analysis, rooted in the inherent variability of handwriting across individuals and writing conditions. Stroke thickness, character shape, spacing irregularities, and background noise all vary substantially from one writer to another, making reliable recognition far harder than its printed-text counterpart.

This thesis presents a feature-enhanced HTR system that addresses both character-level ambiguity and word-level contextual inaccuracy within a unified, end-to-end pipeline. The proposed architecture consists of seven stages: dataset preparation, image preprocessing, CNN-based spatial feature extraction, Bidirectional Long Short-Term Memory (BiLSTM) sequence modeling, hierarchical feature refinement using a Hierarchical Recurrent Neural Network (HRNN), soft attention, and NLP-driven post-processing for contextual error correction.

The system is trained and evaluated on the IAM Handwriting Word Dataset, which provides word images from multiple writers and thus captures a broad range of writing styles. A four-step preprocessing pipeline---grayscale conversion, Gaussian noise filtering, Otsu binarization, and pixel normalization---standardizes the input before feature extraction begins. The CNN encoder extracts spatially discriminative features, which the BiLSTM converts into context-rich sequential representations. HRNN operates across multiple temporal abstraction levels to refine alignment between predicted characters and their spatial positions, while the soft attention module allows the network to focus selectively on informative regions of the feature map. Connectionist Temporal Classification (CTC) decoding is used to produce the final text output without requiring character-level segmentation. An NLP post-processing layer then applies spell correction and language-model-guided context refinement to reduce residual word-level errors.

Experiments show a Character Error Rate (CER) of 0.0102 and a Word Error Rate (WER) of 0.0847 after NLP correction, representing a relative WER reduction of approximately 44\% over the uncorrected output. These results demonstrate that architectural feature enhancement and linguistic post-processing are complementary, and that combining them within a single modular pipeline yields recognition quality that neither component achieves independently.

\vspace{1.0cm}
\noindent \textbf{Keywords:} Handwritten Text Recognition, CNN, BiLSTM, Attention Mechanism, CTC Decoding, NLP Post-Processing, HRNN, IAM Dataset, Deep Learning, Sequence Modeling.

\newpage

% --- 7. LIST OF ABBREVIATIONS ---
\begin{center}
    \Large{\textbf{LIST OF ABBREVIATIONS}} \\
    \vspace{1.0cm}
\end{center}
\setstretch{1.2}
\begin{longtable}{ll}
    \textbf{Abbreviation} & \textbf{Description} \\
    \midrule
    API & Application Programming Interface \\
    BiLSTM & Bidirectional Long Short-Term Memory \\
    CER & Character Error Rate \\
    CNN & Convolutional Neural Network \\
    CRNN & Convolutional Recurrent Neural Network \\
    CTC & Connectionist Temporal Classification \\
    DNN & Deep Neural Network \\
    GRU & Gated Recurrent Unit \\
    HMM & Hidden Markov Model \\
    HRNN & Hierarchical Recurrent Neural Network \\
    HTR & Handwritten Text Recognition \\
    IAM & Institut f{\"u}r Angewandte Mathematik (IAM Handwriting Database) \\
    LSTM & Long Short-Term Memory \\
    MCA & Master of Computer Applications \\
    MITE & Mangalore Institute of Technology \& Engineering \\
    NLP & Natural Language Processing \\
    OCR & Optical Character Recognition \\
    ReLU & Rectified Linear Unit \\
    RNN & Recurrent Neural Network \\
    SVM & Support Vector Machine \\
    VTU & Visvesvaraya Technological University \\
    WER & Word Error Rate \\
\end{longtable}

\newpage

% --- 8. LIST OF SYMBOLS ---
\begin{center}
    \Large{\textbf{LIST OF SYMBOLS}} \\
    \vspace{1.0cm}
\end{center}
\setstretch{1.2}
\begin{longtable}{lll}
    \textbf{Symbol} & \textbf{Description} & \textbf{Unit / Notation} \\
    \midrule
    $\mathbf{x}$ & Input image pixel matrix & Matrix \\
    $\mathbf{W}$ & Weight matrix & Matrix \\
    $\mathbf{b}$ & Bias vector & Vector \\
    $\mathbf{h}_t$ & Hidden state at time step $t$ & Vector \\
    $\alpha_t$ & Attention weight at step $t$ & $[0, 1]$ \\
    $\mathbf{c}_t$ & Context vector at time step $t$ & Vector \\
    $y_t$ & Output prediction at step $t$ & Categorical \\
    $\sigma$ & Sigmoid activation function & Operator \\
    $\tanh$ & Hyperbolic tangent function & Operator \\
    $\mathcal{L}_{CTC}$ & CTC loss function & Scalar \\
    $N$ & Number of training samples & Integer \\
    $T$ & Sequence length & Integer \\
    $C$ & Number of character classes & Integer \\
    $d$ & Edit distance between two strings & Integer \\
\end{longtable}

\newpage

% --- 9. LATEX LISTS AND TABLE OF CONTENTS ---
\setstretch{1.1}
\addtocontents{toc}{\protect\thispagestyle{plain}}
\tableofcontents
\newpage
\addcontentsline{toc}{chapter}{List of Figures}
\listoffigures
\newpage
\addcontentsline{toc}{chapter}{List of Tables}
\listoftables
\newpage
